\documentclass[10pt]{article}
\usepackage[margin=0.78in]{geometry}
\usepackage{amsmath,amssymb,booktabs,array,microtype}
\usepackage[hidelinks]{hyperref}
\usepackage{xcolor}
\newcommand{\statusbox}[1]{\fcolorbox{black}{gray!8}{\parbox{0.94\linewidth}{#1}}}
\title{CertVIC: Separating Visual Responsiveness from Intervention Specificity\\
in Open Vision-Language Models}
\author{Anonymous}
\date{}

\begin{document}
\maketitle

\statusbox{\textbf{Evidence status.} This V11 draft reports a real, pilot-scale set of
open-model outputs, but it is not a submission-ready empirical paper. The item-validity
screen used to form the 91-item intervention set was machine-assisted, the independent
second-rater sheet is blank, and the 30-item Spurious V2 follow-up has no model outputs or
human labels. Numerical confidence-sequence crossings are therefore reported separately
from full scientific claim eligibility.}

\begin{abstract}
Vision-language model evaluations often reward a response change after an image edit,
but responsiveness alone is ambiguous: a model may react to any visible perturbation
rather than to the intended semantic intervention. CertVIC separates two paired
quantities: responsiveness to edits that should change an answer and specificity under
edits that should leave it unchanged. A pilot on three open 7--8B models yields large
descriptive intervention gaps on 91 machine-screened ADE20K pairs, with anytime-valid
lower bounds above the configured numerical threshold. The specificity control changes
the interpretation. On 94 irrelevant-edit pairs, Qwen2.5-VL-7B flips on 12 items
(0.1277) and fails the frozen observed-rate gate of 0.10; InternVL2-8B and
LLaVA-OneVision-7B flip on 1 and 3 items. All twelve Qwen failures are Qwen-only within
the paired three-model comparison, and parser/provenance checks do not explain them.
These results motivate a model-dependent specificity hypothesis, not a broad robustness
claim. We describe the evidence boundary, the prospective control protocol, and the
human and scale validations required before stronger conclusions.
\end{abstract}

\section{Why responsiveness and specificity must be separated}

An edited image can change a model answer for at least two reasons. The intended semantic
fact may have changed, or the model may simply be sensitive to a low-level perturbation.
Conversely, an unchanged answer after a relevant edit may indicate decision inertia even
when the model recognizes natural object absence. A useful intervention evaluation must
therefore measure both responsiveness and invariance under a matched no-change control.

CertVIC's primary contribution is an evaluation protocol and estimand, not a collection
of execution wrappers. The protocol links a source image, a documented edit, a locked
question/answer pair, raw model responses, strict parsing, and a claim gate. Its
statistical component supplies a time-uniform uncertainty bound for a predeclared item
order. The pilot's scientific observation is narrower: responsiveness and specificity
are empirically separable, and specificity is model dependent on the available control.

\paragraph{Scope relative to prior work.}
The repository's related-work matrix identifies counterfactual VQA, visual robustness,
consistency testing, and anytime-valid inference as the closest categories, but verified
bibliographic anchors are still missing. This draft therefore makes no priority claim.
A submission must establish which prior methods measure answer updates without a
matched irrelevant-edit gate, and which uncertainty procedures remain valid under
continuous monitoring.

\section{Paired estimands and statistical decision}

For eligible item $i$, let $A_i$ be original-image correctness and let $C_i$ indicate
whether the parse-valid response pair changes for a relevant intervention or remains
unchanged for an irrelevant control. The historical $C_i$ does not require the edited
answer to equal edited gold; correct semantic update success is reported separately.
The descriptive quantities are
\begin{equation}
 a=\frac{1}{n}\sum_i A_i,\qquad
 p=\frac{1}{n}\sum_i C_i,\qquad
 \Delta=a-p.
\end{equation}
The historical symbol $p$ denotes consistency, not a p-value. For a no-change control,
$F_i=1-C_i$ and $\hat f=n^{-1}\sum_i F_i$ is the spurious flip rate.

The implementation maps the gap variable to
\begin{equation}
 D_i=(A_i-C_i+1)/2\in[0,1],\qquad
 \Delta=2\mathbb{E}[D_i]-1,
\end{equation}
and applies a Hoeffding normal-mixture confidence sequence with a data-independent scale
set by the planned horizon. Its time-uniform coverage supports optional stopping for the
predeclared order under the bounded conditional-mean assumptions. It does not establish
edit validity, use of the intended visual evidence, provider reproducibility, or population
generalization. Those require separate gates.

\paragraph{Frozen V1 rule and prospective V2 rule.}
The V1 decision remains the observed-rate rule $\hat f\leq0.10$; it is not changed after
seeing outcomes. For V2, Qwen is the predeclared primary model and the prospective rule
uses a one-sided 95\% Clopper--Pearson upper bound that must be at most 0.10. A joint
three-model statement uses Bonferroni level $0.05/3$ per model. Missing or unparseable
pairs fail the primary endpoint. Raw and pre-result validity-filtered analyses are both
reported.

\section{Repository evidence and validity status}

The intervention pilot contains 91 item pairs selected from 103 candidates. The source
images and masks are from the local ADE20K release; edits include removal, displacement,
and occlusion. The predictions are real outputs imported from free Kaggle runs for
Qwen2.5-VL-7B, InternVL2-8B, and LLaVA-OneVision-7B. Their exact model commit revisions
were not recorded, which limits replication across changing model repositories.

The repository labels this item set ``human reviewed,'' but its completed review CSV
identifies the reviewer as \texttt{assistant\_visual\_review\_v1}. The independent
second-rater export contains 91 blank rows. V11 therefore treats the model predictions
as real observed evidence and the review-derived inclusion decisions as
machine-assisted preliminary. This distinction weakens claim eligibility without
altering the raw responses or the numerical estimators.

The V1 specificity arm contains 94 paired irrelevant edits per model. All 188 response
rows per provider parse as strict yes/no tokens; item IDs, providers, image pairs, and
merged zip members match. An objective image audit finds no target-mask overlap, but 20
of 94 difference bounding boxes intersect conservative target boxes. The detectability
diagnostic reports symmetric grouped-item AUC 0.7123, below its repository flag threshold
of 0.8 but clearly not
proof of imperceptibility.

\section{Pilot results}

\begin{table}[t]
\centering
\caption{Real pilot outputs. $a$, $p$, and $\Delta$ use the 91 intervention pairs. The
CS lower bound is a numerical time-uniform bound, not full evidence eligibility. The V1
specificity gate is the frozen observed-rate rule.}
\small
\begin{tabular}{lrrrrrrl}
\toprule
Model & $n$ & $a$ & $p$ & $\Delta$ & CS LB & V1 flips & Gate \\
\midrule
Qwen2.5-VL-7B & 91 & .923 & .176 & .747 & .364 & 12/94 & fail \\
InternVL2-8B & 91 & .923 & .099 & .824 & .441 & 1/94 & pass \\
LLaVA-OneVision-7B & 91 & .890 & .176 & .714 & .331 & 3/94 & pass \\
\bottomrule
\end{tabular}
\end{table}

The intervention gaps are large for every model in this pilot, and the numerical
confidence-sequence lower bounds exceed the configured gap threshold 0.05. Those values
are reproducible from the raw paired rows. They do not clear the full policy: the set has
$n=91$ versus the configured minimum 150, lacks completed independent human validation,
and Qwen fails the separate specificity gate.

For Qwen and InternVL, every observed answer change on this arm also moves from the
original gold answer to the edited gold answer (16/91 and 9/91). LLaVA changes on 16/91
items, but only 13 end on the correct edited answer. This distinction is why $p$ is
described as raw response change rather than semantic success.

The natural present/absent control indicates that the models can recognize unedited
absence: combined accuracy is 0.897 for Qwen, 0.935 for InternVL, and 0.932 for LLaVA on
369 validation images. Thus the observed intervention behavior is not explained solely
by an unconditional tendency to answer that the queried object is present. The control
does not by itself validate the edited images.

\paragraph{Paired model comparison.}
On the same 94 specificity items, the Qwen-minus-InternVL flip-risk difference is 0.117
and the Qwen-minus-LLaVA difference is 0.096. Exact McNemar p-values are 0.0034 and
0.0352; after a three-comparison Holm family the second contrast is not below 0.05.
These analyses are retrospective and exploratory. InternVL versus LLaVA has risk
difference $-0.021$ and exact p-value 0.625.

\section{Forensics of the twelve Qwen flips}

Each Qwen failure changes from parsed ``yes'' on the original image to ``no'' on the
edited image. The corresponding InternVL and LLaVA pairs do not flip. No missing row,
duplicate item/variant key, parse error, provider mismatch, or image-path mismatch was
found. Two Qwen failures have conservative difference bounding boxes intersecting the
target box, while target-mask overlap is zero for every item. Machine-generated visual
categories in older artifacts are diagnostics and are not used as human-validity labels.

This pattern is consistent with model-dependent sensitivity to irrelevant edits, but it
does not identify a mechanism. The model revisions, quantization details, and exact
processor commits are incomplete in the historical run manifests. A repeat under pinned
revisions and a blinded validity review is required before assigning the difference to
architecture rather than run configuration.

\section{Prospective Spurious V2 and scale decision}

The repository contains a stricter 30-item V2 package built from the 94-item source pool.
Its objective construction enforces zero target-mask overlap, no target-box intersection,
minimum box distance 75 pixels, and a stored detectability-score cap of 0.12. The set was
constructed after V1 was observed. Four of the twelve known Qwen failures remain in the
set and eight are removed by post-V1 geometry/salience rules. It is therefore a
retrospective stricter-control sensitivity set, not an independent confirmatory sample.

No V2 provider output or real human label is present. Moreover, $n=30$ is too small for a
strong simultaneous three-model upper-bound statement: even zero failures cannot put a
one-sided Bonferroni upper bound below 0.10. V2 is useful as a re-encoding/strict-filter
sensitivity check, but it cannot clear a confirmatory specificity gate. A larger
independently sourced control set containing no V1 items is needed for a broad
specificity claim.

Human review must precede result unblinding. Two independent raters receive randomized
anonymous pairs without provider outputs or V1 failure labels. They judge target
preservation, answer invariance, perturbation acceptability, ambiguity, answerability,
and retention. The raw set remains reported regardless of exclusions.

Main-500 remains blocked. It becomes scientifically useful only after the specificity
branch is resolved, item selection and replacement are outcome blind, model revisions
are pinned, and real human review and objective edit-quality gates pass. A Qwen failure
does not require abandoning the project; it changes the paper's question from uniform
specificity to when responsiveness and specificity diverge across models.

\section{Limitations and research path}

The current evidence has five dominant limitations. First, human item validity is not
complete. Second, the sample is small and from one dataset/domain. Third, the V2 set is
derived from a post-V1 source pool and has only 30 items. Fourth, historical model and
processor revisions are not pinned. Fifth, related work and legal redistribution details
still require verified sources. Diagnostic polarity and mechanism probes also contain
format-specific parse failures and cannot support mechanistic claims.

The shortest credible path is to repair and hash-lock the execution notebook and
transactional importer; complete blinded review before any new unblinding; use the
30-item runs only if their retrospective sensitivity value justifies the compute; build
an independent control set; apply the frozen decision there; and only then decide whether
a locked Main-500 or a small second-domain confirmation adds more value. The current
local assets do not include unused control candidates or a second-domain dataset, so the
independent set requires new data acquisition outside this local-only pass.

\section{Conclusion}

The real pilot shows large descriptive update gaps and model-dependent specificity under
the frozen V1 rule. Stronger conclusions require blinded human review, pinned replication,
and an independent powered control.

\end{document}
